% !TEX program = xelatex
% 五張獨立圖：7×8 格（與 defense_game/grid.py 相同）
%   xelatex defender_movement_diagrams.tex
\documentclass[11pt]{article}
\usepackage[a4paper,landscape,margin=14mm]{geometry}
\usepackage{fontspec}
\usepackage{xeCJK}
\setmainfont{Times New Roman}
\setCJKmainfont{Songti TC}
\setCJKsansfont{Heiti TC}
\usepackage{tikz}
\usetikzlibrary{arrows.meta,calc}

\pagestyle{empty}

% 與遊戲一致：COLS=7（球門短邊）、ROWS=8（向球門縱深）
\newcommand{\pitchcols}{7}
\newcommand{\pitchrows}{8}
\newcommand{\cellsize}{1.12cm}
\newcommand{\pitchmx}{0.75}  % 左／下留空給座標
\newcommand{\pitchmy}{0.55}

% 格 (x,y) 中心座標（須在 pitch scope 內使用）
\newcommand{\gc}[2]{#1+0.5,#2+0.5}

\newcommand{\defenderpage}[4]{%
  \clearpage
  \noindent
  {\fontsize{22}{26}\selectfont\bfseries #1\par}
  \vspace{0.15em}
  {\fontsize{14}{18}\selectfont #2\par}
  \vspace{0.45em}
  \begin{center}
  \begin{tikzpicture}[
    x=\cellsize, y=\cellsize,
    every node/.style={font=\small}
  ]
    \begin{scope}[shift={(\pitchmx,\pitchmy)}]
      #3
    \end{scope}
    \foreach \x in {0,...,6} {
      \node[font=\scriptsize, anchor=north]
        at (\pitchmx+\x+0.5, \pitchmy-0.32) {\x};
    }
    \foreach \y in {0,...,7} {
      \node[font=\scriptsize, anchor=east]
        at (\pitchmx-0.32, \pitchmy+\y+0.5) {\strut\y};
    }
  \end{tikzpicture}
  \end{center}
  \vspace{0.4em}
  \noindent{\normalsize #4\par}
}

% 7 欄 × 8 列：格 (x,y) = 矩形 [x,x+1]×[y,y+1]，與 grid.py 一致
\newcommand{\pitchgrid}{%
  \foreach \x in {1,...,6} {
    \draw[gray!35, line width=0.35pt] (\x,0) -- (\x,\pitchrows);
  }
  \foreach \y in {1,...,7} {
    \draw[gray!35, line width=0.35pt] (0,\y) -- (\pitchcols,\y);
  }
  \draw[line width=1pt] (0,0) rectangle (\pitchcols,\pitchrows);
  % 球門三格頂邊加粗（不填灰、不畫出界括號）
  \draw[line width=2.4pt] (2,\pitchrows) -- (5,\pitchrows);
  \node[font=\footnotesize, anchor=south west]
    at (5.08,\pitchrows+0.02) {球門 $(2,7)(3,7)(4,7)$};
}

\newcommand{\deftoken}[3]{%
  \draw[thick, fill=white] (\gc{#1}{#2}) circle (0.24);
  \node[font=\normalsize\bfseries] at (\gc{#1}{#2}) {#3};
}

% 進攻球員：白底圓＋圈內密斜線（不標名字／座標；持球也用同符號，不加圓心黑點）
\newcommand{\atktoken}[2]{%
  \coordinate (atk@) at (#1+0.5,#2+0.5);
  \begin{scope}
    \clip (atk@) circle (0.24);
    \fill[white] (atk@) circle (0.24);
    \foreach \i in {-6,-5,...,6} {
      \draw[line width=0.35pt, black!75]
        ($(atk@)+({-0.36+0.075*\i},-0.36)$)
        -- ($(atk@)+({0.36+0.075*\i},0.36)$);
    }
  \end{scope}
  \draw[thick] (atk@) circle (0.24);
}
\newcommand{\ballatk}[2]{\atktoken{#1}{#2}}

% 箭頭：每段只跨 1 格（切比雪夫距離 1）；可選中間點標記第 2 步
\newcommand{\moveone}[4]{%
  \draw[-{Latex}, line width=1.5pt] (\gc{#1}{#2}) -- (\gc{#3}{#4});
}
\newcommand{\movetwo}[6]{%
  \draw[-{Latex}, line width=1.5pt] (\gc{#1}{#2}) -- (\gc{#3}{#4}) -- (\gc{#5}{#6});
  \fill (\gc{#3}{#4}) circle (0.07);
}

\begin{document}
\thispagestyle{empty}
\noindent{\fontsize{26}{30}\selectfont\bfseries 防守角色移動示意（共五張）\par}
\vspace{0.3em}
\noindent{\normalsize
\textbf{$7\times 8$ 格}（寬 $7$、深 $8$；短邊為球門，與解防守題程式相同）。
$y$ 愈大愈近球門；座標與棋盤一致。箭頭每跨\textbf{一邊}為 1 格（逼搶／遠距攔截可兩段，中點為第 2 步）。觸發：結束回合／傳球／高空傳球後。\par}

% ========== 1 Presser ==========
\defenderpage{1. 瘋狗逼搶員}{Presser — 每回合最多 2 格}{
  \pitchgrid
  \ballatk{4}{3}
  \draw[dashed, thick] (0,4) -- (\pitchcols,4);
  \node[anchor=west, font=\footnotesize, fill=white, inner sep=1pt, align=left]
    at (4.85,5.15) {不可使 $y$ 大於持球者\\（界線 $y{=}4$）};
  \deftoken{1}{3}{P}
  \movetwo{1}{3}{2}{3}{3}{3}
  \deftoken{3}{3}{P}
  \node[font=\footnotesize, anchor=east, fill=white, inner sep=1pt]
    at ($(3.5,3.5)+(-0.38,0)$) {鄰格施壓};
}{朝持球者靠近，\textbf{不踩持球格}；$y$ 不得大於持球者。}

% ========== 2 Block ==========
\defenderpage{2. 區域大閘}{Block — 每回合橫移 1 格}{
  \pitchgrid
  \draw[thick, gray!60] (0,5) -- (\pitchcols,5);
  \node[anchor=east, font=\footnotesize] at (-0.05,5.5) {固定橫線 $y{=}5$};
  \ballatk{5}{6}
  \draw[dashed, thick] (\gc{5}{6}) -- (\gc{3}{7});
  \node[font=\footnotesize, anchor=west, fill=white, inner sep=1pt]
    at ($(3.5,7.5)+(0.35,0)$) {球門中心};
  \deftoken{3}{5}{B}
  \moveone{3}{5}{4}{5}
  \deftoken{4}{5}{B}
}{只沿\textbf{同一橫線}滑動，站到「球→球門中心」走廊與該線交點。定點大閘不移動。}

% ========== 3 Shadow ==========
\defenderpage{3. 影子盯人}{Shadow — 每回合 1 格}{
  \pitchgrid
  \atktoken{4}{2}
  \draw[dashed, thick] (\gc{4}{2}) -- (\gc{3}{7});
  \deftoken{2}{2}{S}
  \moveone{2}{2}{3}{3}
  \deftoken{3}{3}{S}
  \node[font=\footnotesize, anchor=north] at (\gc{4}{1}) {門側格；不可站到較小 $y$};
}{朝被盯者與球門之間的\textbf{門側格}移動；不站進攻者本人所在格。}

% ========== 4 Interceptor ==========
\defenderpage{4. 路線攔截者}{Interceptor — 遠 2 格／近 1 格}{
  \pitchgrid
  \ballatk{1}{2}
  \atktoken{5}{5}
  \draw[dashed, thick] (\gc{1}{2}) -- (\gc{5}{5});
  \deftoken{1}{2}{I}
  \movetwo{1}{2}{2}{2}{3}{3}
  \deftoken{3}{3}{I}
  \node[font=\footnotesize, anchor=west, fill=white, inner sep=1pt]
    at ($(3.5,3.5)+(0.32,0.15)$) {傳球線中點};
}{鎖定最危險前鋒，移向「持球者—該隊友」地滾傳球線的\textbf{中點}。}

% ========== 5 Goalkeeper ==========
\defenderpage{5. 守門員}{Goalkeeper — 每回合橫移 1 格}{
  \pitchgrid
  \ballatk{1}{2}
  \deftoken{4}{7}{G}
  \moveone{4}{7}{3}{7}
  \deftoken{3}{7}{G}
}{僅在球門三格內，沿球門線跟隨球的\textbf{左右}位置（遠近皆追）。}

\end{document}
